\par\Needspace{6\baselineskip}\medskip
\SourceLeaf{142}\pdfbookmark[2]{102.}{article-102}
\begin{center}102.\end{center}
Da $y$ bei der Hyperbel ebenso wie bei der Ellipse positiv sein muss, ist die Lösung der Gleichung 16 auch hier eindeutig\footnote{Es braucht kaum erwähnt zu werden, dass unsere Tafel II bei der Hyperbel ebenso wie bei der Ellipse zur Lösung dieser Gleichung verwendet werden kann, solange $h$ innerhalb ihrer Grenzen bleibt.}. Bei Gleichung $16^*$ ist hier jedoch etwas anders zu argumentieren als bei der Ellipse. Aus der Theorie der Gleichungen lässt sich leicht zeigen, dass diese Gleichung bei positivem $H$\HistoricalNoteDoubleUnbracketed{Die Größe $H$ kann offenbar nur negativ werden, wenn $\zeta>\tfrac16$ ist. Einem solchen Wert von $\zeta$ entspräche ein Wert von $z$, der größer als $2{,}684$ wäre und somit weit außerhalb der Grenzen dieses Verfahrens läge.}, sofern sie überhaupt eine positive reelle Wurzel hat, eine negative und zwei positive Wurzeln besitzt. Die beiden positiven Wurzeln sind entweder gleich, und zwar beide $=\tfrac16\BareRoot5-\tfrac16=0{,}20601$, oder die eine liegt über und die andere unter dieser Grenze. Wir zeigen nun, dass bei unserer Aufgabe stets die größere Wurzel zu nehmen ist. Der Aufgabe liegt die Annahme zugrunde, dass $z$ klein und jedenfalls nicht größer als $0{,}3$ ist, damit die dritte Tafel verwendbar bleibt. Setzt man in Gleichung $13^*$ für $M$ den Ausdruck $Y\BareRoot{(L+z)}$ ein, so erhält man
\[
Y+1=(L+z)Z>(1+z)Z,
\]
\noindent oder
\[
Y>\tfrac13-\frac4{3\mathpunct{.}5}z+\frac{4\mathpunct{.}6}{3\mathpunct{.}5\mathpunct{.}7}zz-\frac{4\mathpunct{.}6\mathpunct{.}8}{3\mathpunct{.}5\mathpunct{.}7\mathpunct{.}9}z^3+\text{usw.},
\]
\noindent woraus leicht folgt, dass für so kleine Werte von $z$, wie wir sie hier voraussetzen, stets $Y>0{,}20601$ sein muss. Tatsächlich ergibt die Rechnung, dass für die Gleichheit von $(1+z)Z$ mit dieser Grenze $z=0{,}79858$ sein müsste. So große Werte von $z$ liegen jedoch weit außerhalb des Bereichs, für den unser Verfahren gedacht ist.
